\chapter*{Conclusion Générale}
\addcontentsline{toc}{chapter}{Conclusion Générale}
\markboth{Conclusion Générale}{Conclusion Générale}

Ce rapport a proposé une exploration structurée et progressive du thème de l'exploitation des données à l'aide des modèles d'intelligence artificielle, depuis les fondements théoriques jusqu'à la concrétisation dans la plateforme \textbf{AgriYield}, en passant par la présentation des grandes familles de modèles mobilisables et la méthodologie opérationnelle.

\bigskip
Le \textbf{premier chapitre} a posé les bases conceptuelles indispensables~: la notion de donnée sous ses différentes formes, les exigences de qualité qui conditionnent la fiabilité des analyses, et le processus général d'exploitation qui transforme les données brutes en connaissances actionnables. Il a mis en évidence les limites structurelles des approches classiques face aux caractéristiques des données contemporaines (les cinq V du Big Data), justifiant ainsi le recours aux modèles d'intelligence artificielle.

\bigskip
Le \textbf{deuxième chapitre} a dressé un panorama des principaux modèles d'IA, en distinguant les approches supervisées (régression, classification), non supervisées (clustering, réduction de dimensionnalité) et les architectures de deep learning. Il a montré que chaque famille présente des forces et des limites qui rendent le choix du modèle indissociable de la compréhension fine du problème et des données.

\bigskip
Le \textbf{troisième chapitre} a décrit une méthodologie rigoureuse et reproductible, couvrant la formulation du problème, la collecte et la compréhension des données, le prétraitement, l'analyse exploratoire, l'entraînement des modèles, l'évaluation des performances et l'interprétation des résultats. La méthodologie CRISP-DM a servi de cadre structurant pour cette démarche itérative.

\bigskip
Le \textbf{quatrième chapitre} a concrétisé l'ensemble de ces concepts dans l'application AgriYield~: un système de prédiction du rendement agricole entraîné sur le jeu de données FAO/Kaggle couvrant 101 pays sur 23 ans. Le modèle XGBoost sélectionné atteint un coefficient de détermination $R^2 = 0.91$, confirmant la pertinence de l'approche. Les résultats ont permis d'identifier les facteurs climatiques (pluviométrie, température) comme les variables les plus déterminantes du rendement, et d'alimenter des recommandations concrètes pour l'allocation des intrants et la planification des politiques agricoles.

\bigskip
Le \textbf{cinquième chapitre} a ouvert une réflexion critique sur les apports, les limites (granularité des données, concept drift, enjeux éthiques) et les perspectives d'amélioration, notamment l'intégration de données satellites en temps réel, le déploiement MLOps et les outils d'explicabilité.

\bigskip
\noindent À travers cette démarche, ce rapport confirme que la science des données, soutenue par les modèles d'intelligence artificielle, offre un cadre puissant pour transformer des données brutes en leviers de décision à fort impact sociétal. Cependant, la qualité des résultats dépend fondamentalement de la \textbf{rigueur méthodologique}, de la \textbf{qualité des données} mobilisées et de la \textbf{capacité à interpréter les résultats} dans leur contexte. Ces trois dimensions constituent le c\oe{}ur de compétence du praticien en science des données, et elles ont guidé chaque étape du travail présenté dans ce rapport.

\bigskip
L'agriculture camerounaise et africaine plus largement dispose d'un potentiel considérable que les outils d'IA, bien appliqués, peuvent contribuer à révéler et à optimiser~: c'est l'ambition portée par la plateforme AgriYield.

\newpage
